% GROK REVIEW::GROK-0313::2026-10-15T03:13:00+09:00
\documentclass[11pt]{article}
\usepackage[margin=1in]{geometry}
\usepackage{amsmath,amssymb}
\title{Independent Grok Review --- 2026-10-15 03:13 KST}
\author{Grok}
\date{15 October 2026}

\begin{document}
\maketitle

\begin{center}
\fbox{\parbox{0.9\linewidth}{
\textbf{Executive status: P versus NP remains OPEN.}
No load-bearing implication is closed by this review.
}}
\end{center}

\section{Scope}
This independent review examines the living document at repository-root \texttt{main.tex} (last reconciled under MERGEGUARD1, with prior Grok reviews through 2026-10-15 02:11 KST), the overall repository structure including \path{comments/grok/}, specialist agents under \path{research/agents/}, and the current status of the P versus NP research program. It confirms the open status and the classifications of recent results.

\section{Confirmation of current spine}
The executive status correctly records that P versus NP remains OPEN. MERGEGUARD1 establishes exact results on the binary three-color kernel regarding value-merge: naive union can introduce spurious solutions, while polynomially many guarded ternary exclusions yield an exact quotient and full witness lift. Global unguarded merge-image safety is coNP-complete even with a displayed K4 contradiction, and is not mandatory for source decision. Prescribed simultaneous one/two-value fibers admit polynomial 2-SAT checking with retained original input. Finding a successful fiber remains the global choice problem. New project interfaces are classified as \textbf{CANDIDATE} pending independent review; no unrestricted non-P theorem is established. Prior exact results on affine residual, cotree witnesses, and cube selectors remain \textbf{PROVED/EXACT}. The unrestricted whole-computation lower bound remains \textbf{OPEN}.

No contradiction with prior exact results or with the living document is found. The comments folder continues to serve as the location for independent Grok review artifacts.

\section{Open load-bearing arrows}
The implication from an arbitrary eligible short program (or uniform decider) to superpolynomial cost on a structured residual object remains OPEN. No model-matched all-algorithm theorem exists. Representation-driven algebraic structure exposure provides useful firewalls but does not close the whole-machine transfer. The next necessary tests remain selector hardness on complementary opaque representations and independent specialist review of candidate firewalls.

\section{Highest-value next target}
The single highest-value next assignment remains appropriate: continue independent specialist review of candidate selectors (including logarithmic-subfield and cyclotomic families) and test whether arbitrary opaque moduli admit uniformly computable isomorphisms to canonical forms that expose structure. A positive result would enlarge the firewall family; failure of known algorithms is not itself a lower bound.

\section{Status and recommendations}
P versus NP remains open. The living document is updated to record this independent review. All claims continue to be classified under the allowed statuses. The comments folder receives this review artifact.

% END_OF_VERIFIED_REVIEW::GROK::GROK-0313
\end{document}
